\documentclass[letterpaper]{article}

% -- Packages --
\usepackage[utf8]{inputenc}
\usepackage[T1]{fontenc}
\usepackage{microtype}
\usepackage{amsmath,amssymb,amsfonts,amsthm}
\usepackage{mathtools}
\usepackage{bm}
\usepackage{graphicx}
\usepackage{subcaption}
\usepackage{booktabs}
\usepackage{multirow}
\usepackage{xcolor}
\usepackage{hyperref}
\usepackage{cleveref}
\usepackage{natbib}
\usepackage{url}
\usepackage{enumitem}

% AAAI style (use official template when available; fallback to geometry)
\usepackage[margin=1in]{geometry}

% -- Macros --
\newcommand{\xx}{\mathbf{x}}
\newcommand{\uu}{\mathbf{u}}
\newcommand{\ff}{\mathbf{f}}
\newcommand{\norm}[1]{\left\|#1\right\|}
\newtheorem{empirical}{Empirical Law}
\newtheorem{remark}{Remark}

% -- Title / Author --
\title{How Many Steps? A Universal, Noise-Governed Law\\for Test-Time Compute in Iterative Denoisers}

\author{%
  Hafez Khatib \\
  American University of Beirut \\
  \texttt{hma164@mail.aub.edu}
}

\begin{document}
\maketitle

\begin{abstract}
How many refinement steps should an iterative denoiser take?
We show the optimal depth follows a universal power law
$K^{*}(\sigma) \approx C\sigma^{\alpha}$.
Across four datasets (CIFAR-10, FashionMNIST, MNIST, CelebA-64),
four energy-model head types (KAN, GELU, SiLU, Tanh),
and a non-energy score network, we find $\alpha \approx 1.37$
with cross-architecture variance of just $0.2\%$.
The exponent is universal not only across architectures but across
datasets: CIFAR-10 $\alpha{=}1.377$, FashionMNIST $\alpha{=}1.372$,
MNIST $\alpha{=}1.356$ --- despite very different spectral content
($\hat\beta$ ranges from $2.28$ to $2.90$).
When we probe a pretrained DDPM's score function (35.7M params)
with our iteration protocol, it gives $\alpha{=}1.366$
($R^2{=}0.994$), matching our 35K-param models across a
$1000\times$ scale gap.
At higher resolution (CelebA-64), $\alpha$ shifts to $1.265$,
suggesting weak resolution dependence.
The law holds for all stochastic corruptions
(Gaussian, speckle, Poisson; $R^2 > 0.95$) and persists when
blur is combined with noise, but breaks for purely deterministic
operators (blur-only, JPEG).
Since $\alpha$ is near-constant, a practitioner needs only the
prefactor $C$ --- estimable from a single noise level ---
to set the optimal inference budget at any $\sigma$.
This one-shot calibration matches the accuracy of a full
seven-point fit ($6.4\%$ vs.\ $6.1\%$ mean K* error)
while eliminating the multi-point sweep entirely.
\end{abstract}

%==========================================================================
\section{Introduction}
\label{sec:intro}
%==========================================================================

Test-time compute scaling --- allocating additional computation at
inference rather than training --- has dramatically improved LLM
quality \citep{snell2024llm}, yet vision models remain largely
one-shot: once trained, quality is capped by model capacity.
Energy-based models (EBMs) \citep{lecun2006ebm} run gradient descent
on a learned energy $E_\theta(\xx)$ at inference, offering a natural
mechanism for quality to improve with step count $K$.
Score-based models \citep{song2019generative} and diffusion models
\citep{ho2020ddpm} share the same iterative structure.
Recently, \citet{ma2025inference} and \citet{singhal2025fk}
have explored inference-time scaling for diffusion models,
and \citet{tong2025ld3} learn discretization schedules.
\textbf{None characterize a closed-form power law
governing the optimal step count.}

\textbf{This paper asks:} Is there a predictive,
architecture-independent scaling law that governs the optimal
inference depth for iterative denoisers?

We report an affirmative answer: the optimal number of gradient
steps follows $K^*(\sigma) \approx C\sigma^\alpha$ with a
\emph{near-universal} exponent $\alpha \approx 1.37$.
This universality extends across architectures, model families
(energy-based and score-based), and natural-image datasets at
matched resolution.
Only resolution appears to shift $\alpha$ (32px: $\sim$1.37,
64px: $\sim$1.27).

\textbf{Contributions.}
\begin{enumerate}[leftmargin=*,topsep=2pt,itemsep=1pt,parsep=0pt]
\item \textbf{A universal inference scaling law}
  $K^{*}(\sigma) \approx C\sigma^{\alpha}$ with
  $\alpha = 1.377 \pm 0.003$ across four EBM activation variants
  (KAN, GELU, SiLU, Tanh) and one score network,
  three seeds each on CIFAR-10, confirmed on FashionMNIST
  ($\alpha{=}1.372$) and MNIST ($\alpha{=}1.356$).

\item \textbf{Cross-family and cross-scale unification:} the same
  exponent holds for a direct score network (no scalar potential)
  and for a pretrained DDPM ($1000\times$ larger),
  proving the law transcends both parameterization and model scale.

\item \textbf{A corruption phase diagram:} the law holds for
  all stochastic corruptions (including blur+noise at $R^2 > 0.97$)
  and breaks only for purely deterministic operators.

\item \textbf{Honest falsification} of the spectral prediction
  $\alpha = 2\gamma/\beta$: it fails for FashionMNIST (25\% error)
  and on synthetic GRF $\beta$-sweeps (slope $\approx$ 0,
  expected $-1$). We present an interpretive connection to
  Landweber regularization theory instead.

\item \textbf{A one-shot calibration rule:} fix
  $\alpha{=}1.37$, fit $C$ from a single noise level.
  This matches the 7-point oracle ($6.4\%$ vs.\ $6.1\%$ mean
  $K^*$ error) while being $7\times$ cheaper.
\end{enumerate}

%==========================================================================
\section{Related Work}
\label{sec:related}
%==========================================================================

\paragraph{Test-time compute scaling.}
In language models, additional inference compute via search or
verifiers can match the gains from scaling parameters
\citep{snell2024llm}.
In vision, \citet{ma2025inference} explore noise-conditioned search
for diffusion models and \citet{singhal2025fk} propose FK steering.
IterRef \citep{iterref2025} applies iterative refinement to
discrete diffusion.
None characterize a power law.

\paragraph{Step scheduling for diffusion.}
LD3 \citep{tong2025ld3} learns discretization schedules for
pretrained diffusion models.
\citet{yuan2026cvpr} propose instance-aware discretization.
\citet{pei2025steps} analyze step allocation in flow matching.
Our work differs: we learn the law from the energy/score landscape
itself, without auxiliary training.

\paragraph{Energy-based models.}
EBMs have a long history \citep{lecun2006ebm,du2019implicit}.
\citet{nijkamp2019anatomy} study short-run MCMC for EBM sampling.
\citet{goyal2025ebt} introduce Energy-Based Transformers.
None connect EBM inference depth to test-time scaling laws.

\paragraph{Classical regularization.}
Landweber iteration on linear inverse problems has optimal stopping
$K^* \propto \sigma^{-2/(2s+1)}$ where $s$ is a source-condition
smoothness parameter \citep{engl1996regularization,raskutti2014early}.
We show this classical power-law form holds precisely in modern
nonlinear architectures.

%==========================================================================
\section{Setup and Protocol}
\label{sec:setup}
%==========================================================================

\paragraph{Iterative denoising.}
Given a noisy image $\uu_0 = \xx + \sigma\epsilon$,
we iterate gradient descent on the learned energy:
$\uu_{t+1} = \uu_t - \eta_t \nabla E(\uu_t)$
with $\eta_t = 0.05 \cdot 0.97^t$ (geometric decay),
for up to $K_{\max} = 30$ steps.
For score models, the update is $\uu_{t+1} = \uu_t + \eta_t s_\theta(\uu_t)$.

\paragraph{Optimal depth $K^*(\sigma)$.}
For each test image and noise level $\sigma$, we record the
per-image step that maximizes PSNR: $K^*_i = \arg\max_k \text{PSNR}(\uu_k, \xx_i)$.
The reported $K^*(\sigma)$ is the mean over images and evaluation
noise seeds: $K^*(\sigma) = \frac{1}{N\cdot S}\sum_{i,s} K^*_{i,s}$.
This \emph{per-image continuous} $K^*$ --- as opposed to the argmax
of the mean PSNR curve --- enables genuine per-seed variance estimation.

\paragraph{Fitting.}
We evaluate at $\sigma \in \{0.05, 0.08, 0.12, 0.16, 0.20, 0.25, 0.30\}$
and fit $\log K^* = \alpha \log\sigma + \log C$ via OLS.

\paragraph{Training.}
All models use denoising score matching (DSM) at a single noise level
$\sigma{=}0.15$, AdamW optimizer, cosine LR schedule, 40 epochs,
batch size 32.
Architectures share a $16$-filter, $5{\times}5$ conv backbone
followed by per-pixel heads of varying type.

\paragraph{Architectures tested.}
\begin{itemize}[leftmargin=*,topsep=2pt,itemsep=1pt]
\item \textbf{KAN-EBM}: B-spline KAN head $[48, 16]$ (110K params)
\item \textbf{ConvMLP-GELU/SiLU/Tanh}: MLP head with specified activation (35K params)
\item \textbf{ScoreNet}: Direct score prediction, no scalar energy (77K params)
\end{itemize}

%==========================================================================
\section{Results}
\label{sec:results}
%==========================================================================

\subsection{Architecture Independence (CIFAR-10)}
\label{sec:arch}

We train four EBM architectures and one score network, each with
three independent seeds, under the identical protocol.
\Cref{tab:tier0} shows the results.

\begin{table}[t]
\centering\small
\begin{tabular}{lrr}
\toprule
Architecture & $\alpha$ (mean $\pm$ std) & $R^2$ \\
\midrule
KAN-EBM (110K)     & $1.376 \pm 0.003$ & $0.981$ \\
ConvMLP-GELU (35K) & $1.378 \pm 0.003$ & $0.981$ \\
ConvMLP-SiLU (35K) & $1.379 \pm 0.003$ & $0.981$ \\
ConvMLP-Tanh (35K) & $1.373 \pm 0.003$ & $0.981$ \\
\midrule
ScoreNet (77K)     & $1.376 \pm 0.002$ & $0.982$ \\
\midrule
\textbf{Grand mean} & $\bm{1.377 \pm 0.003}$ & \\
\bottomrule
\end{tabular}
\caption{\textbf{Architecture independence.} Four EBM heads and one
score network, each with 3 seeds, on CIFAR-10 (32$\times$32).
All $\alpha$ values agree within noise ($\Delta\alpha < 0.006$).}
\label{tab:tier0}
\end{table}

\textbf{Cross-family validation.}
The score network has no scalar energy and no gradient structure ---
it directly predicts the noise residual.
Its exponent ($\alpha = 1.376 \pm 0.002$) is statistically
identical to the energy-based family, confirming the law transcends
the energy-based parameterization.

\subsection{Dataset Universality and Resolution Dependence}
\label{sec:datasets}

We train ConvMLP-GELU and ScoreNet on three additional datasets
(FashionMNIST, MNIST at 28$\times$28; CelebA at 64$\times$64),
each with 3 seeds.
\Cref{tab:datasets} shows $\alpha$ is near-universal across
datasets at comparable resolution.

\begin{table}[t]
\centering\small
\begin{tabular}{lrrr}
\toprule
Dataset & Resolution & $\alpha$ (mean $\pm$ std) & $\hat\beta$ \\
\midrule
CIFAR-10     & $32\times32$ & $1.377 \pm 0.003$ & 2.90 \\
FashionMNIST & $28\times28$ & $1.372 \pm 0.006$ & 2.28 \\
MNIST        & $28\times28$ & $1.356 \pm 0.008$ & 2.87 \\
\midrule
CelebA-64    & $64\times64$ & $1.265 \pm 0.007$ & 3.46 \\
\bottomrule
\end{tabular}
\caption{\textbf{Dataset universality.}
$\alpha$ is near-constant ($\sim$1.37) across three datasets at
small resolution, despite very different spectral exponents $\hat\beta$.
Only the higher-resolution CelebA-64 deviates ($\alpha = 1.265$).}
\label{tab:datasets}
\end{table}

\textbf{Key observation:}
FashionMNIST has $\hat\beta = 2.28$ (much lower than CIFAR's $2.90$),
yet gives $\alpha = 1.372$ --- virtually identical to CIFAR's $1.377$.
If $\alpha$ were spectrally determined (as our earlier theory predicted),
Fashion should have $\alpha \approx 1.8$.
This falsifies the spectral prediction $\alpha = 2\gamma/\beta$
(\Cref{sec:mechanism}).

\textbf{Resolution dependence.}
CelebA-64 at $64\times64$ gives $\alpha = 1.265$, significantly lower
than the $\sim$1.37 cluster.
Architecture independence is preserved at this resolution
(score: $1.257$, ConvMLP: $1.268$, KAN: $1.271$).
This suggests $\alpha$ has weak resolution dependence ---
an open question for future work.

\subsection{Corruption Phase Diagram}
\label{sec:phase}

We test the law under seven corruption types.
\Cref{tab:phase} summarizes the results.

\begin{table}[t]
\centering\small
\begin{tabular}{lrrr}
\toprule
Corruption & $\alpha$ & $R^2$ & Verdict \\
\midrule
Gaussian noise     & 1.47 & 0.994 & \textsc{Law} \\
Speckle noise      & 1.26 & 0.981 & \textsc{Law} \\
Poisson noise      & 1.59 & 0.991 & \textsc{Law} \\
\midrule
Blur + noise (all levels) & --- & $>$0.97 & \textsc{Law} \\
\midrule
Blur only          & --- & $<$0.60 & Breaks \\
JPEG only          & --- & $<$0.10 & Breaks \\
\bottomrule
\end{tabular}
\caption{\textbf{Corruption phase diagram.}
The law holds for all stochastic corruptions and for blur+noise
composites. It breaks only for purely deterministic operators.}
\label{tab:phase}
\end{table}

\textbf{Blur + noise surprise.}
The law survives even substantial blur ($\sigma_{\text{blur}} = 3.0$)
combined with additive noise ($R^2 > 0.97$ for all three model types
at all blur levels).
The stochastic noise component dominates.
Simple operational rule: if the corruption has a stochastic additive
component, the law holds.

\subsection{Scale Transfer: Pretrained DDPM (35.7M params)}
\label{sec:ddpm}

The strongest test of universality is whether the law holds on a
model $1000\times$ larger than our training-controlled models.
We evaluate \texttt{google/ddpm-cifar10-32} (35.7M params,
\citealt{ho2020ddpm}) under our exact K-sweep protocol.

\textbf{Important clarification:}
We do \emph{not} use the DDPM's native ancestral sampler or DDIM.
Instead, we extract the DDPM's learned score function
$s_\theta(\uu, t)$ and use it as the gradient oracle in our
standard iteration: $\uu_{k+1} = \uu_k + \eta_k s_\theta(\uu_k, t)$
with geometric step decay.
This tests whether the DDPM's learned score landscape obeys the
same power law when probed with our protocol, not whether the
DDPM's native sampling schedule follows the law.

We test two conditioning protocols:
\begin{itemize}[leftmargin=*,topsep=2pt,itemsep=1pt]
\item \textbf{V1 (fixed conditioning):} condition at $t{=}42$
  ($\sigma \approx 0.15$ in the DDPM noise schedule), matching
  our EBMs' single-$\sigma$ training.
\item \textbf{V2 (noise-matched):} condition at the diffusion timestep
  whose scheduled noise level matches each evaluation $\sigma$.
\end{itemize}

\begin{table}[t]
\centering\small
\begin{tabular}{lrrr}
\toprule
Protocol & $\alpha$ & $R^2$ & Params \\
\midrule
Our EBMs (mean, 4 arch) & $1.377 \pm 0.003$ & $0.981$ & 35--110K \\
DDPM V1 (fixed cond.)   & $1.366$ & $0.994$ & 35.7M \\
DDPM V2 (noise-matched) & $1.205$ & $0.995$ & 35.7M \\
\bottomrule
\end{tabular}
\caption{\textbf{Scale transfer.}
A pretrained DDPM ($1000\times$ our model size) gives
$\alpha = 1.37$ under fixed conditioning ---
matching our small models exactly.
Noise-matched conditioning shifts $\alpha$ to $1.21$.}
\label{tab:ddpm}
\end{table}

\textbf{V1 matches our small models exactly:}
$\alpha_{\text{DDPM-V1}} = 1.366$ vs.\
$\alpha_{\text{EBM}} = 1.377 \pm 0.003$ ---
within $0.8\%$.
The power law transfers across $1000\times$ parameter scale,
from 35K-param energy models to a 35.7M-param pretrained
diffusion model, with $R^2 = 0.994$.

\textbf{Conditioning matters:}
V2 (noise-matched) gives $\alpha = 1.205$, significantly lower.
When the model's conditioning adapts to each noise level,
the effective iteration dynamics change.
This is analogous to the resolution effect (CelebA-64
$\alpha = 1.27$): both represent changes to what the model
``sees'' at each step.

\subsection{Feedforward Contrast}
\label{sec:dncnn}

A DnCNN-style feedforward denoiser (65K and 189K params) produces
comparable or better PSNR at $K{=}1$ but has no $K$-dial: it cannot
trade compute for quality.
The scaling law is specific to iterative architectures that admit
a refinement loop.

\subsection{One-Shot Calibration Rule}
\label{sec:oneshot}

Since $\alpha \approx 1.37$ is near-constant across architectures
and datasets, the calibration simplifies dramatically:
fix $\alpha = 1.37$ and estimate only $C$ from a \emph{single}
noise level $\sigma_{\text{cal}}$:
\begin{equation}
C = K^*(\sigma_{\text{cal}}) / \sigma_{\text{cal}}^{1.37}
\end{equation}
then predict $K^*(\sigma) = C\sigma^{1.37}$ for any $\sigma$.

\Cref{tab:calibration} compares calibration methods across 24 runs
(12 CIFAR-10 + 6 FashionMNIST + 6 MNIST).

\begin{table}[t]
\centering\small
\begin{tabular}{lrrr}
\toprule
Method & \#$\sigma$ needed & Mean $K^*$ err & Max $K^*$ err \\
\midrule
One-shot ($\alpha{=}1.37$, fit $C$)  & 1 & 6.4\% & 22.8\% \\
7-pt oracle (fit $\alpha + C$)       & 7 & 6.1\% & 25.4\% \\
Fixed $K{=}5$                        & 0 & 126\% & 471\%  \\
\bottomrule
\end{tabular}
\caption{\textbf{Calibration budget comparison.}
The one-shot rule (1 noise level) matches the full oracle
(7 noise levels) because the power-law fit itself has
$\sim$20\% residual error at extreme $\sigma$ values.
More calibration points do not help.}
\label{tab:calibration}
\end{table}

\textbf{Key insight:} the one-shot rule matches the oracle because
the power law has $\sim$20--25\% residual error at extreme $\sigma$
regardless of fitting method.
The practical gain is $7\times$ fewer calibration evaluations,
with no loss in prediction accuracy.

Per-dataset breakdown: CIFAR-10 mean error $8.5\%$,
FashionMNIST $5.8\%$, MNIST $2.8\%$.

%==========================================================================
\section{Why $\alpha \approx 1.37$? Connection to Classical Regularization}
\label{sec:mechanism}
%==========================================================================

\paragraph{Interpretive framing (not a derivation).}
In classical Landweber iteration for linear inverse problems
$A\xx = \mathbf{b} + \sigma\epsilon$, optimal early stopping gives
$K^* \propto \sigma^{-2/(2s+1)}$, where $s$ is the source condition
smoothness \citep{engl1996regularization,raskutti2014early}.
There, $K^*$ \emph{decreases} with $\sigma$: more noise means
stop earlier to prevent noise amplification.

Our setting is different.
We start from the noisy observation $\uu_0 = \xx + \sigma\epsilon$
and run gradient descent on a learned energy $E(\uu)$, with
geometrically decaying step sizes $\eta_t = \eta_0 \rho^t$.
Here $K^*$ \emph{increases} with $\sigma$: more noise means
greater initial displacement from the energy minimum, requiring
more gradient steps to reach it.
The effective regularization in our scheme comes from the
\emph{accumulated displacement} $\sum_{t=0}^{K-1} \eta_t
= \eta_0(1-\rho^K)/(1-\rho)$, which increases with $K$.
Defining the residual regularization as
$\lambda(K) \propto 1/\sum\eta_t$, we have $\lambda$ decreasing
with $K$.
The classical optimality condition
$\lambda^* \propto \sigma^{2/(2s+1)}$ then requires $K^*$ to
\emph{increase} with $\sigma$ (to drive $\lambda$ down),
recovering the correct sign.
For $\alpha \approx 1.37$, the implied source condition is
$s \approx 0.23$ --- a weak regularity consistent with
natural images in the relevant Sobolev scale.

\paragraph{Why is $\alpha$ universal?}
The near-constancy of $\alpha$ across datasets (despite different
$\hat\beta$) suggests the effective source condition $s$ is
a property of the \textbf{training recipe and iteration scheme}
(DSM + geometric-decay gradient descent), not the data's spectral
structure.
All models share the same learning algorithm and decay schedule,
which may impose a universal implicit regularization.

\paragraph{Falsified spectral prediction.}
Our earlier theory predicted $\alpha = 2\gamma/\beta$, where
$\gamma$ is the learned curvature spectrum exponent and $\beta$ is
the data power-spectrum exponent.
This theory makes two testable predictions, both of which fail:

\begin{enumerate}[leftmargin=*,topsep=2pt,itemsep=1pt]
\item \textbf{Cross-dataset:} FashionMNIST has $\hat\beta = 2.28$
  (vs.\ CIFAR's $2.90$), predicting $\alpha_{\text{Fashion}} \approx 1.8$.
  Measured: $1.372$ (25\% error).

\item \textbf{Synthetic $\beta$-sweep:} On Gaussian random fields with
  controlled $\beta$, the measured slope of $\log\alpha$ vs.\ $\log\beta$
  is $-0.09$ (expected: $-1.0$).
\end{enumerate}

We report these negative results transparently.
The near-universality of $\alpha$ is more fundamental than
spectral scaling, and its origin remains an open question.

\paragraph{Open question.}
Why does resolution shift $\alpha$ ($32$px: $1.37$, $64$px: $1.27$)?
One hypothesis: higher resolution admits more informative high-frequency
modes, effectively raising $s$ and lowering $\alpha$.

%==========================================================================
\section{Limitations}
\label{sec:limitations}
%==========================================================================

\textbf{Scale.}
Our training-controlled models are 30--110K parameters on 28--64px images.
A pretrained 35.7M-param DDPM confirms the law at $1000\times$ scale,
but verification on even larger models (EDM, DiT) at higher resolution
remains future work.

\textbf{Single-$\sigma$ training.}
Our standard recipe trains at $\sigma = 0.15$.
Multi-$\sigma$ training changes $\alpha$ (preliminary: $\alpha \approx 1.7$
for rich recipes); full characterization is future work.

\textbf{Noise estimation.}
The practical rule assumes $\sigma$ is known.
Blind noise estimation adds uncertainty not modeled here.

\textbf{Theoretical mechanism.}
The Landweber connection is interpretive, not a derivation.
The spectral prediction $\alpha = 2\gamma/\beta$ is falsified.
The origin of the universal $\alpha \approx 1.37$ remains open.

%==========================================================================
\section{Conclusion}
\label{sec:conclusion}
%==========================================================================

We characterize a universal inference-depth scaling law
$K^*(\sigma) \approx C\sigma^\alpha$ for iterative denoisers,
with $\alpha \approx 1.37$ near-constant across four EBM head types,
a score network, four datasets, and a pretrained DDPM whose score
function --- probed with our iteration protocol --- gives the same
exponent across a $1000\times$ parameter-scale gap.
The exponent is not set by the data's spectral structure ---
our spectral prediction theory is falsified --- but appears to be
an emergent property of the training and iteration scheme.
The law provides a practical one-shot calibration rule:
fix $\alpha = 1.37$, measure $C$ from one noise level,
and predict the optimal step count for any noise level.
This connects the active research area of test-time compute scaling
to classical regularization theory, and delivers a deployable tool
for practitioners using iterative denoisers.

\bibliographystyle{abbrvnat}
\bibliography{references}

\end{document}
